\begin{enumerate}

\item \textbf{Generalisation to modern targeted-alpha-therapy isotopes
      (Ac-225, Th-227).}\\
      \emph{Basis:} The paper's RBE=14.7 is a single-point measurement
      at Am-241's 130~keV/$\mu$m LET; clinical TAT isotopes deliver
      alphas at $\sim$55--70~keV/$\mu$m, closer to the LET-RBE peak.
      Whether the SW-1573 clonogenic RBE translates to Ac-225-PSMA-617
      or Th-227 conjugates is untested.\\
      \emph{Next steps:} Meta-analyse published SW-1573 (or matched
      squamous carcinoma) survival data under Ac-225 and Th-227, reweight
      via ICRU-40 or Katz LET-RBE curves, compare predicted vs measured
      survival. Free endpoint: literature + LQ arithmetic.

\item \textbf{Reconciliation with mechanistic DNA-damage-repair models
      (MEDRAS, TOPAS-nBio + BioDose, McMahon RMF).}\\
      \emph{Basis:} The 14.7-fold RBE gap between DSB induction (RBE=1)
      and cell killing (RBE=14.7) is exactly what modern mechanistic
      repair-misrepair-fixation models are designed to predict.\\
      \emph{Next steps:} Run open-source MEDRAS at 130~keV/$\mu$m,
      compare its predicted ``complex DSB fraction'' to the paper's
      empirical $\sim$10$\times$ lethal-DSB enrichment, publish
      reconciliation. Free endpoint: MEDRAS on laptop CPU.

\item \textbf{Sensitivity of the inferred $\beta_\gamma$ (and $\alpha/\beta
      \approx 1.57$~Gy) to the assumed iso-effect level.}\\
      \emph{Basis:} Fig-2's ``RBE=4'' does not specify an effect level.
      Alternative conventions ($D_{37}$, $D_{50}$, iso-2Gy) yield
      different $\beta_\gamma$ values and different $\alpha/\beta$ ratios.\\
      \emph{Next steps:} Recompute under each convention, produce a
      sensitivity table, cross-check against SW-1573 $\alpha/\beta$ from
      independent literature. Free endpoint: extends existing
      \texttt{code/pass2\_extended\_claims.py}.

\item \textbf{Cell-line dependence: is SW-1573 typical or an outlier?}\\
      \emph{Basis:} RBE varies by factors of 2--3 across cell lines for
      the same alpha source. SW-1573 alone cannot represent lung-cancer
      alpha radiobiology; TAT dose planning uses tissue-averaged RBE.\\
      \emph{Next steps:} Mine open-access alpha-clonogenic papers on
      A549, H460, H1299, SK-MES-1 at 100--150~keV/$\mu$m; extract or
      digitise their LQ fits; compare RBE-at-10\%-survival distribution
      to SW-1573's 14.7. Free endpoint: literature review + LQ
      arithmetic.

\item \textbf{ML surrogate for LQ parameters from cell-line features.}\\
      \emph{Basis:} A Gaussian-process or random-forest regressor trained
      on $\sim$30 published high-LET LQ studies (features: p53 status,
      ploidy, doubling time, initial DSB yield) could interpolate the
      paper's fit to unmeasured cell lines and LETs without new
      experiments.\\
      \emph{Next steps:} Curate high-LET LQ dataset from open literature,
      encode cell-line features from Cellosaurus (free API), train GP
      with LET covariate, predict SW-1573 fit and compare. Free endpoint:
      scikit-learn on laptop CPU.

\end{enumerate}
